\documentclass{article}
\usepackage{mgvalues}\renewcommand{\thetable}{S\arabic{table}}

\begin{document}

\section{Supplementary methods}

\subsection{Table S1. Code lists used to identify the outcome}

Everything under \texttt{manuscript/supplementary/} is read by the same gates as the paper ---
a number here is a binding or a defect, exactly as in \texttt{main.tex} --- and built as its
own document, so it does not count against the journal's word limit and can be uploaded to
the supplementary slot rather than pasted onto the end of the manuscript.

The case definition belongs here rather than in the Methods because it is a list, and because
a reader checking it wants the whole list rather than a sentence about it. It is emitted from
the terms the analysis filtered on, so the definition in this table and the definition the
code applied cannot come apart.

\begin{table}\caption{Code lists used to identify the outcome.}\label{tab:codes}\mg{table.outcome_codes}\end{table}

Events were coded with the data generator's own list of event terms, release
\mg{results.dictionary.generator_event_terms}; a real study codes with MedDRA, and names the
version the events were coded with.

\subsection{Table S2. The adjusted model}

The model as it was fitted: the outcome, and each variable with the kind it was declared and
how the model entered it. The table is made from the fit and from those declarations, not
typed, so it cannot describe a model other than the one that produced the estimate.

\begin{table}\caption{The adjusted model.}\label{tab:model}\mg{table.model_adjusted}\end{table}

\subsection{Sensitivity of the estimate to the interval quoted}

The primary analysis reports a 95\% confidence interval of \mg{results.ror.ci_low} to
\mg{results.ror.ci_high}. The narrower 90\% confidence interval, \mg{results.ror.ci90_low} to
\mg{results.ror.ci90_high}, is given for readers comparing this estimate with signal-detection
thresholds that use it. Both exclude the null, so the conclusion does not turn on which is
quoted.

% Each interval is one interval() call and carries its level, so quoting either one upper
% bound first is refused. A comment, because this paragraph is about the toolkit rather than
% the study, and its own percentages, 95% and 90%, would otherwise be reported.

\subsection{Reporting}

This supplement is referred to from the Methods and is listed in the submission manifest as
its own file, alongside the manuscript rather than inside it.

\end{document}
